\documentclass[12pt,a4paper]{report}

% ===================== ENCODAGE ET LANGUE =====================
\usepackage[utf8]{inputenc}
\usepackage[T1]{fontenc}
\usepackage[english,french]{babel} % francais = langue principale (dernier charge)

% ===================== MISE EN PAGE =====================
\usepackage[a4paper,margin=2.5cm]{geometry}
\usepackage{setspace}
\setstretch{1.15}

% ===================== POLICE (Arial) =====================
% Helvetica (helvet) est metriquement equivalente a Arial ; scaled=1 pour
% garder la taille nominale (pas de reduction automatique de l'x-height).
\usepackage[scaled=1]{helvet}
\renewcommand{\familydefault}{\sfdefault}

% ===================== TAILLE DES TITRES =====================
% Template : 14pt pour les titres (chapitres, sections), 12pt pour les
% sous-titres (sous-sections et niveaux inferieurs).
\usepackage{titlesec}
\titleformat{\chapter}[hang]{\normalfont\bfseries\fontsize{14}{16.8}\selectfont}{\thechapter.}{10pt}{}
\titleformat{\section}{\normalfont\bfseries\fontsize{14}{16.8}\selectfont}{\thesection}{8pt}{}
\titleformat{\subsection}{\normalfont\bfseries\fontsize{12}{14.4}\selectfont}{\thesubsection}{6pt}{}
\titleformat{\subsubsection}{\normalfont\bfseries\fontsize{12}{14.4}\selectfont}{\thesubsubsection}{6pt}{}
% \paragraph et \subparagraph sont utilises dans ce rapport comme des
% intitules autonomes (suivis d'un saut de ligne, parfois d'un tableau ou
% d'une figure) plutot que du texte encha\^in\'e : on les garde donc en
% [hang], comme les autres niveaux, plut\^ot qu'en [runin].
\titleformat{\paragraph}[hang]{\normalfont\bfseries\fontsize{12}{14.4}\selectfont}{\theparagraph}{6pt}{}
\titleformat{\subparagraph}[hang]{\normalfont\bfseries\itshape\fontsize{12}{14.4}\selectfont}{\thesubparagraph}{6pt}{}

% ===================== GRAPHIQUES ET COULEURS =====================
\usepackage{graphicx}
\graphicspath{{images/}}
\usepackage{xcolor}
% Identit\'e visuelle "Le Sceau" (m\^eme palette que la plateforme CyberGuard
% B\'enin) : papier chaud, bleu marine, accent or.
\definecolor{brandnavy}{HTML}{16304A}
\definecolor{brandgold}{HTML}{B8863A}
\definecolor{brandpaper}{HTML}{FFFDF7}
\definecolor{brandborder}{HTML}{DCD3B8}
\usepackage{tikz}
\usetikzlibrary{shapes,arrows.meta,positioning,calc}

% ===================== TABLEAUX =====================
\usepackage{booktabs}
\usepackage{longtable}
\usepackage{array}

% ===================== LISTES =====================
\usepackage{enumitem}

% ===================== CODE SOURCE =====================
\usepackage{listings}
\lstdefinestyle{code}{
  basicstyle=\ttfamily\footnotesize,
  keywordstyle=\color{blue}\bfseries,
  commentstyle=\color{gray}\itshape,
  stringstyle=\color{purple!70!black},
  numbers=left,
  numberstyle=\tiny\color{gray},
  frame=single,
  breaklines=true,
  breakatwhitespace=false,
  showstringspaces=false,
  tabsize=2,
  backgroundcolor=\color{black!3}
}
\lstset{style=code}

% ===================== EN-TETES / PIEDS DE PAGE =====================
\usepackage{fancyhdr}
\pagestyle{fancy}
\fancyhf{}
\fancyhead[L]{\small\leftmark}
\fancyfoot[C]{\thepage}
\renewcommand{\headrulewidth}{0.4pt}

% ===================== FIGURES ET FLOTTANTS =====================
\usepackage{float}
\usepackage{caption}
\captionsetup{font=small,labelfont=bf}

% ===================== LIENS =====================
\usepackage{hyperref}
\hypersetup{
  colorlinks=true,
  linkcolor=black,
  citecolor=black,
  urlcolor=blue,
  pdftitle={CyberGuard Benin - Rapport de stage},
  pdfauthor={Orace ATTINMAWASSONOU, Fructueux HOUNGBEME}
}

% ===================== INFORMATIONS DU RAPPORT =====================
% Point unique de modification : change ici, ca se repercute partout.
\newcommand{\etudiantA}{Orace ATTINMAWASSONOU}
\newcommand{\etudiantB}{Fructueux HOUNGBEME}
\newcommand{\ecole}{Ecole Nationale d'Economie Appliquee et du Management (ENEAM)}
\newcommand{\universite}{Universite d'Abomey Calavi (UAC)}
\newcommand{\filiere}{Informatique de gestion}
\newcommand{\anneeacademique}{2025-2026}
\newcommand{\theme}{CyberGuard B\'enin \\ Plateforme intelligente d'analyse et d'\'evaluation de la s\'ecurit\'e des sites web}
\newcommand{\entreprise}{Direction G\'en\'erale de la Recherche Scientifique et de l'Innovation (DGRSI)}
\newcommand{\maitredestage}{KOUTON Joslyn}
\newcommand{\datedebut}{06 juillet 2026}
\newcommand{\datefin}{31 juillet 2026}

% ===================== FIGURE PLACEHOLDER =====================
% Insere la vraie image si le fichier existe deja dans images/, sinon
% affiche un encadre "a inserer" -- permet d'ajouter les images au fur et
% a mesure sans jamais retoucher les chapitres, et le document continue de
% compiler proprement sur Overleaf tant que certaines manquent encore.
\newcommand{\TODOFigure}[3]{%
  \begin{figure}[H]
    \centering
    \IfFileExists{images/#1}{%
      \includegraphics[width=0.85\textwidth]{#1}%
    }{%
      \fbox{%
        \begin{minipage}{0.85\textwidth}
          \centering
          \vspace{1.5cm}
          \textit{[Image \`a ins\'erer ici : #1]}
          \vspace{1.5cm}
        \end{minipage}%
      }%
    }
    \caption{#2}
    \label{#3}
  \end{figure}
}

\begin{document}

% ===================== PAGE DE GARDE =====================
\begin{titlepage}
\centering

\begin{tikzpicture}[remember picture, overlay]
  \draw[brandborder, line width=0.8pt]
    ($(current page.north west) + (1.2cm,-1.2cm)$) rectangle
    ($(current page.south east) + (-1.2cm,1.2cm)$);
  % Petits losanges dores aux quatre coins du cadre, pour un rendu plus
  % soigne qu'un simple rectangle.
  \draw[brandgold, line width=0.8pt, rotate around={45:($(current page.north west)+(1.2cm,-1.2cm)$)}]
    ($(current page.north west) + (1.2cm,-1.2cm)$) +(-0.13,-0.13) rectangle +(0.13,0.13);
  \draw[brandgold, line width=0.8pt, rotate around={45:($(current page.north east)+(-1.2cm,-1.2cm)$)}]
    ($(current page.north east) + (-1.2cm,-1.2cm)$) +(-0.13,-0.13) rectangle +(0.13,0.13);
  \draw[brandgold, line width=0.8pt, rotate around={45:($(current page.south west)+(1.2cm,1.2cm)$)}]
    ($(current page.south west) + (1.2cm,1.2cm)$) +(-0.13,-0.13) rectangle +(0.13,0.13);
  \draw[brandgold, line width=0.8pt, rotate around={45:($(current page.south east)+(-1.2cm,1.2cm)$)}]
    ($(current page.south east) + (-1.2cm,1.2cm)$) +(-0.13,-0.13) rectangle +(0.13,0.13);
\end{tikzpicture}

\vspace*{1.2cm}

% ===== LOGOS UAC / ENEAM =====
% Depose logo-uac.jpg et logo-eneam.jpg dans images/ : ils s'affichent
% automatiquement a la bonne taille. En attendant, un encadre discret
% indique leur emplacement prevu, sans casser la compilation.
\noindent
\begin{minipage}{0.92\textwidth}
  \centering
  \begin{minipage}[c][2.4cm][c]{4cm}
    \centering
    \IfFileExists{images/logo-uac.jpg}{%
      \includegraphics[height=2.2cm]{logo-uac.jpg}%
    }{%
      \fbox{\begin{minipage}[c][1.8cm][c]{3.2cm}\centering\footnotesize\itshape\color{brandnavy!70} [Logo UAC]\end{minipage}}%
    }
  \end{minipage}%
  \hfill
  \begin{minipage}[c][2.4cm][c]{4cm}
    \centering
    \IfFileExists{images/logo-eneam.jpg}{%
      \includegraphics[height=2.2cm]{logo-eneam.jpg}%
    }{%
      \fbox{\begin{minipage}[c][1.8cm][c]{3.2cm}\centering\footnotesize\itshape\color{brandnavy!70} [Logo ENEAM]\end{minipage}}%
    }
  \end{minipage}
\end{minipage}

\vspace{1cm}

% ===== SCEAU DECORATIF =====
% Petit emblème rond (double liseré or, silhouette de bouclier) en clin
% d'oeil au nom de la charte graphique "Le Sceau" et au theme cybersecurite
% du rapport -- purement ornemental.
\begin{tikzpicture}
  \draw[brandgold, line width=0.7pt] (0,0) circle (0.85cm);
  \draw[brandgold, line width=0.4pt] (0,0) circle (0.72cm);
  \fill[brandnavy] (-0.32,0.4) -- (0.32,0.4) -- (0.32,0.02) -- (0,-0.42) -- (-0.32,0.02) -- cycle;
  \fill[brandgold] (0,0.16) circle (0.07cm);
  \fill[brandgold] (-0.06,0.02) -- (0.06,0.02) -- (0.045,-0.14) -- (-0.045,-0.14) -- cycle;
\end{tikzpicture}

\vspace{0.6cm}

{\color{brandnavy}\footnotesize\bfseries R\'EPUBLIQUE DU B\'ENIN}\\[0.35cm]
{\color{brandgold}\rule{2.6cm}{0.6pt}}\\[0.7cm]

{\color{brandnavy}\bfseries MINIST\`ERE DE L'ENSEIGNEMENT SUP\'ERIEUR ET DE LA\\
RECHERCHE SCIENTIFIQUE}\\[1cm]

{\color{brandnavy}\rmfamily\Large\bfseries \universite}\\[0.35cm]
{\color{brandnavy}\rmfamily\bfseries \ecole}\\[0.4cm]
{\color{brandnavy}\underline{\textbf{Fili\`ere}} : \filiere}\\[1.6cm]

{\color{brandnavy}\footnotesize\bfseries THEME}\\[0.5cm]

\begin{tikzpicture}
  % Rectangle decale en fond (brandborder) pour simuler une legere ombre
  % portee, sans dependre de la librairie tikz "shadows".
  \node[fill=brandborder, sharp corners, inner sep=14pt, text width=0.75\textwidth, xshift=0.12cm, yshift=-0.12cm] {
    \rmfamily\Large\bfseries\color{brandborder} \theme
  };
  \node[draw=brandnavy, line width=0.8pt, sharp corners, fill=brandpaper, inner sep=14pt, text width=0.75\textwidth, align=center] {
    \rmfamily\Large\bfseries\color{brandnavy} \theme
  };
\end{tikzpicture}

\vspace{1.8cm}

{\color{brandgold}\rule{4cm}{0.6pt}}

\vspace{1.8cm}

\begin{minipage}{0.5\textwidth}
  \centering
  {\color{brandnavy}\itshape R\'ealis\'e par~:}\\[0.35cm]
  {\color{brandnavy}\rmfamily\bfseries \etudiantA}\\
  {\color{brandnavy}\rmfamily\bfseries \etudiantB}
\end{minipage}

\vspace{1.6cm}

{\color{brandnavy}\itshape Ma\^itre de stage~:}\\[0.25cm]
{\color{brandnavy}\rmfamily\bfseries \maitredestage}

\vfill

{\color{brandgold}\rule{2.6cm}{0.6pt}}\\[0.4cm]
{\color{brandnavy}\bfseries Ann\'ee acad\'emique~: \anneeacademique}

\vspace{1.2cm}

\end{titlepage}

% ===================== SOMMAIRE (bref) =====================
\chapter*{Sommaire}
\addcontentsline{toc}{chapter}{Sommaire}

\begin{description}[leftmargin=0cm,style=nextline]
  \item[R\'ESUM\'E]
  \item[ABSTRACT]
  \item[INTRODUCTION]
  \item[CHAPITRE I~: D\'EROULEMENT DU STAGE]
  \item[CHAPITRE II~: PROJET DE PROGRAMMATION]
  \item[CHAPITRE III~: R\'ESULTATS OBTENUS]
  \item[CONCLUSION ET PERSPECTIVES]
  \item[R\'EF\'ERENCES BIBLIOGRAPHIQUES]
  \item[ANNEXES]
  \item[TABLE DES MATI\`ERES]
\end{description}


% ===================== DEDICACES =====================
\chapter*{D\'edicaces}
\addcontentsline{toc}{chapter}{D\'edicaces}

Je d\'edie ce modeste travail~:

\`A mon p\`ere, \textbf{ATTINMAWASSONOU S\'eraphin}, pour son soutien ind\'efectible et ses
encouragements tout au long de mon parcours~;

\`A ma m\`ere, \textbf{DONOUVOSSI Nad\`ege}, pour son amour, ses sacrifices et
ses pri\`eres~;

\`A mes s\oe urs, pour leur affection et leurs
encouragements constants~;

\`A mes amis, pour leur soutien moral et leur pr\'esence~;

\`A tous ceux qui, de pr\`es ou de loin, ont contribu\'e \`a la r\'eussite de
ce travail.

\vspace{0.5cm}
\begin{center}
  $\diamond$ $\diamond$ $\diamond$
\end{center}
\vspace{0.5cm}

Je d\'edie ce modeste travail~:

\`A mon p\`ere, \textbf{HOUNGBEME Luc}, pour son soutien ind\'efectible et ses
encouragements tout au long de mon parcours~;

\`A ma m\`ere, \textbf{ZOUNON Fifonsi}, pour son amour, ses sacrifices et
ses pri\`eres~;

\`A mes fr\`eres et s\oe urs, pour leur affection et leurs
encouragements constants~;

\`A mes amis, pour leur soutien moral et leur pr\'esence~;

\`A tous ceux qui, de pr\`es ou de loin, ont contribu\'e \`a la r\'eussite de
ce travail.

% ===================== REMERCIEMENTS =====================
\chapter*{Remerciements}
\addcontentsline{toc}{chapter}{Remerciements}

Nous tenons \`a exprimer notre gratitude \`a~:

\begin{itemize}
  \item Monsieur le \textbf{Directeur de la Recherche Scientifique et de l'Innovation}, pour nous avoir permis de r\'ealiser ce stage au sein de \entreprise{} et pour ses conseils avis\'es~;
  \item \textbf{\maitredestage}, notre ma\^itre de stage, pour son accompagnement et sa disponibilit\'e constante tout au long de ce stage~;
  \item Toute l'\'equipe p\'edagogique de l'\ecole{} et les enseignants de la fili\`ere \filiere, en particulier \textbf{Monsieur COMBLAN Maurice}, chef de d\'epartement~;
  \item Tout le personnel de \entreprise{} pour son accueil, les conseils et le partage de comp\'etences techniques~;
  \item Tous ceux qui ont particip\'e \`a la r\'ealisation de ce travail.
\end{itemize}

% ===================== LISTE DES FIGURES / TABLEAUX =====================
\listoffigures
\listoftables

% ===================== SIGLES ET ABREVIATIONS =====================
\chapter*{Sigles et abr\'eviations}
\addcontentsline{toc}{chapter}{Sigles et abr\'eviations}

\begin{description}[leftmargin=2.5cm,style=multiline]
  \item[API] Application Programming Interface
  \item[CSS] Cascading Style Sheets
  \item[ENEAM] Ecole Nationale d'Economie Appliqu\'ee et de Management
  \item[HTML] HyperText Markup Language
  \item[HTTP] HyperText Transfer Protocol
  \item[HTTPS] HyperText Transfer Protocol Secure
  \item[JSON] JavaScript Object Notation
  \item[ORM] Object-Relational Mapping
  \item[PDF] Portable Document Format
  \item[PHP] Hypertext Preprocessor
  \item[REST] Representational State Transfer
  \item[SGBD] Syst\`eme de Gestion de Base de Donn\'ees
  \item[SPA] Single Page Application
  \item[SQL] Structured Query Language
  \item[SSL/TLS] Secure Sockets Layer / Transport Layer Security
  \item[UAC] Universit\'e d'Abomey Calavi
  \item[UML] Unified Modeling Language
  \item[URL] Uniform Resource Locator
\end{description}


% ===================== RESUME / ABSTRACT =====================
\chapter*{R\'esum\'e}
\addcontentsline{toc}{chapter}{R\'esum\'e}

Au B\'enin, de plus en plus d'entreprises, \'ecoles, administrations et PME se dotent
d'un site web sans disposer des comp\'etences ou des moyens n\'ecessaires pour
v\'erifier si celui-ci respecte les standards de s\'ecurit\'e de base. Les outils existants
(Mozilla Observatory, SSL Labs) sont souvent en anglais, dispers\'es, et peu adapt\'es
\`a un utilisateur non-technicien. Dans le cadre de notre stage acad\'emique, nous avons
con\c{c}u et d\'evelopp\'e \textbf{CyberGuard B\'enin}, une plateforme web qui automatise
le diagnostic de s\'ecurit\'e d'un site : l'utilisateur saisit une URL, et le syst\`eme
ex\'ecute une s\'erie de v\'erifications non intrusives (HTTPS, certificat SSL, en-t\^etes
de s\'ecurit\'e, cookies, exposition d'informations serveur), puis traduit ces r\'esultats
en un score sur 100, un niveau de risque et des recommandations compr\'ehensibles.
Le projet a d\'ebut\'e par une clarification du besoin et des objectifs, suivie d'une
phase de mod\'elisation UML (acteurs, cas d'utilisation, diagrammes de classes,
d'objets et de s\'equence) puis de la conception de la base de donn\'ees. La
r\'ealisation technique repose sur une API Laravel (PHP) c\^ot\'e backend et une
application React c\^ot\'e frontend. Ce projet propose ainsi une solution concr\`ete
pour d\'emocratiser l'acc\`es \`a un premier niveau d'audit de s\'ecurit\'e pour les
propri\'etaires de sites web au B\'enin.

\bigskip
\noindent\textbf{Mots-cl\'es~:} s\'ecurit\'e web, audit automatis\'e, HTTPS, Laravel,
React, UML.

\begin{otherlanguage}{english}
\chapter*{Abstract}
\addcontentsline{toc}{chapter}{Abstract}

In Benin, a growing number of businesses, schools, administrations and SMEs run a
website without the skills or resources needed to verify whether it meets basic
security standards. Existing tools (Mozilla Observatory, SSL Labs) are often
English-only, scattered, and not designed for non-technical users. As part of our
academic internship, we designed and developed \textbf{CyberGuard B\'enin}, a web
platform that automates this security diagnostic: the user enters a URL, and the
system runs a series of non-intrusive checks (HTTPS, SSL certificate, security
headers, cookies, server information disclosure), then translates these results
into a score out of 100, a risk level, and clear, actionable recommendations. The
project began with a clarification of the need and objectives, followed by a UML
modeling phase (actors, use cases, class, object and sequence diagrams) and
database design. The technical implementation relies on a Laravel (PHP) API on
the backend and a React application on the frontend. This project offers a
concrete solution to democratize access to a first level of security auditing for
website owners in Benin.

\bigskip
\noindent\textbf{Keywords:} web security, automated audit, HTTPS, Laravel, React,
UML.
\end{otherlanguage}


% ===================== INTRODUCTION =====================
\chapter*{Introduction}
\addcontentsline{toc}{chapter}{Introduction}
\markboth{Introduction}{}

La formation pratique constitue un pilier fondamental dans le parcours acad\'emique
des \'etudiants, permettant l'acquisition de comp\'etences professionnelles et
l'application concr\`ete des connaissances th\'eoriques acquises en cours. Les stages
repr\'esentent ainsi un pont essentiel entre le monde acad\'emique et le monde
professionnel.

Au B\'enin, de plus en plus d'entreprises, \'ecoles, administrations et PME se dotent
d'un site web pour exister en ligne -- vitrine, services, communication. Mais la
s\'ecurit\'e informatique reste un domaine technique r\'eserv\'e \`a des sp\'ecialistes, et
la majorit\'e de ces structures n'ont ni les comp\'etences ni les moyens de v\'erifier
si leur site respecte les standards de s\'ecurit\'e de base. Un site web mal s\'ecuris\'e
expose son propri\'etaire, et ses visiteurs, \`a des risques concrets~: usurpation via
un site sans HTTPS, interception de donn\'ees sur une connexion non chiffr\'ee,
attaques facilit\'ees par l'absence d'en-t\^etes de s\'ecurit\'e, ou fuite d'informations
sur le serveur utilis\'e. Le probl\`eme n'est pas seulement technique~: c'est un
probl\`eme d'acc\`es \`a l'information. Le propri\'etaire d'un site ne sait g\'en\'eralement
pas si son site est vuln\'erable, ni quoi v\'erifier, ni comment corriger. Les outils
existants (Mozilla Observatory, SSL Labs) sont souvent en anglais, dispers\'es, et
pas pens\'es pour un utilisateur non-technicien qui veut simplement comprendre
\og~est-ce que mon site est s\^ur, et que dois-je faire~?~\fg{}

C'est dans ce contexte que nous avons con\c{c}u \textbf{CyberGuard B\'enin}, une
plateforme qui automatise ce diagnostic. L'utilisateur cr\'ee un compte, saisit
l'URL de son site, et la plateforme lance une s\'erie de v\'erifications non
intrusives (HTTPS, certificat SSL, en-t\^etes de s\'ecurit\'e, exposition
d'informations serveur, cookies). Le syst\`eme traduit ces r\'esultats techniques en
un score sur 100 et un niveau de risque (Bon / Moyen / Faible), accompagn\'e de
recommandations concr\`etes et compr\'ehensibles.

Ainsi, dans le cadre de notre formation, nous avons effectu\'e un stage d'un
mois au sein de la \entreprise{}, du \datedebut{} au \datefin{}, qui nous a
permis de r\'ealiser ce projet. Ce stage s'inscrit dans notre cursus acad\'emique
et constitue une \'etape cruciale de notre parcours d'apprenants.

Affect\'es \`a la Cellule de la Statistique, de la Documentation et du
Pr\'earchivage de la DGRSI, nous avons men\'e en parall\`ele de nos activit\'es
courantes le pr\'esent projet de programmation, CyberGuard B\'enin, dans le cadre
de notre m\'emoire de fin de formation en Licence Professionnelle.

Le pr\'esent rapport, structur\'e en trois parties principales, a pour objectif de
pr\'esenter le travail effectu\'e durant notre stage. Apr\`es avoir pr\'esent\'e le
d\'eroulement du stage, incluant une description de l'entreprise d'accueil, de son
environnement technologique et des missions qui nous ont \'et\'e confi\'ees, nous
abordons la description et l'analyse conceptuelle qui ont permis l'aboutissement
du projet~: l'\'etude des besoins, l'analyse fonctionnelle, ainsi que la
mod\'elisation du syst\`eme \`a travers les divers diagrammes UML (cas d'utilisation,
classes, objets, s\'equence). Enfin, nous pr\'esentons les \'etapes de
l'impl\'ementation proprement dite, les choix technologiques effectu\'es, les
difficult\'es rencontr\'ees et les solutions apport\'ees, ainsi que les r\'esultats
obtenus et les perspectives d'am\'elioration.


% ===================== CHAPITRES =====================
\chapter{D\'eroulement du stage}
\label{chap:deroulement}

Durant un mois, du \datedebut{} au \datefin{}, notre stage au sein de la
\entreprise{} nous a permis de d\'ecouvrir le fonctionnement d'une structure
publique en charge du pilotage de la recherche scientifique au B\'enin, tout en
menant en parall\`ele notre projet de m\'emoire de fin de formation. Ce chapitre
pr\'esente en premier lieu le minist\`ere de tutelle et la DGRSI, puis d\'ecrit les
activit\'es men\'ees durant le stage. Il se termine par une analyse des apports de
cette exp\'erience et des difficult\'es rencontr\'ees.

\section{Pr\'esentation de la structure}

\subsection{Historique}

La Direction G\'en\'erale de la Recherche Scientifique et de l'Innovation
(DGRSI) est l'une des deux directions g\'en\'erales du Minist\`ere de
l'Enseignement Sup\'erieur et de la Recherche Scientifique (MESRS), aux c\^ot\'es
de la Direction G\'en\'erale de l'Enseignement Sup\'erieur (DGES). Son
organisation et son fonctionnement actuels sont r\'egis par le d\'ecret
n\textsuperscript{o}2023-150 du 12 avril 2023 portant attributions,
organisation et fonctionnement du MESRS, qui pr\'ecise les missions confi\'ees \`a
chacune de ses structures.

\subsection{Raison d'\^etre}

Le MESRS a pour mission de concevoir, mettre en \oe uvre, suivre et \'evaluer la
politique nationale en mati\`ere d'enseignement sup\'erieur, de recherche
scientifique et d'innovation, conform\'ement aux lois en vigueur en R\'epublique
du B\'enin. Au sein de ce minist\`ere, la DGRSI est sp\'ecifiquement charg\'ee de la
conception, de la coordination et du suivi de la mise en \oe uvre de la
politique de l'\'Etat en mati\`ere de recherche scientifique et d'innovation.

\subsection{Objectifs}

Parmi les attributions confi\'ees \`a la DGRSI figurent notamment~: l'\'elaboration
de la politique nationale de la recherche et de l'innovation, en collaboration
avec l'Agence B\'eninoise pour la Recherche et l'Innovation (ABRI)~; la
coordination et la supervision des activit\'es de recherche men\'ees dans les
institutions universitaires publiques et priv\'ees ainsi que dans les instituts,
centres et laboratoires de recherche~; le suivi de ces activit\'es aux plans
global et sectoriel~; la promotion de la formation en sciences et techniques,
notamment \`a l'endroit des femmes et des jeunes~; et l'\'elaboration d'une
plateforme nationale d'informations en science, technologie et innovation.

\section{Organisation technique et manag\'eriale}

\subsection{Structure organisationnelle}

Outre le secr\'etariat de la direction, la DGRSI s'appuie sur des structures
d'accompagnement (Cellule informatique~; Cellule de la Statistique, de la
Documentation et du Pr\'earchivage~; assistant-r\'egisseur) et deux d\'epartements
techniques~: le D\'epartement de la Qualit\'e et du Suivi de la Recherche
scientifique et de l'Innovation, et le D\'epartement du Partenariat et de la
Promotion de la Recherche Scientifique et de l'Innovation.

Nous avons \'et\'e affect\'es \`a la \textbf{Cellule de la Statistique, de la
Documentation et du Pr\'earchivage}, charg\'ee notamment d'\'elaborer les fiches
statistiques sur les activit\'es de la direction, en vue de la collecte et du
traitement des informations, de tenir \`a jour et de pr\'earchiver les documents
de la direction, de produire les statistiques du secteur et de veiller \`a leur
archivage.

\TODOFigure{organigramme_MESRS_DGRSI.png}{Organigramme simplifi\'e du MESRS et de la DGRSI -- la case fonc\'ee indique la structure d'accueil du stage (source~: d\'ecret n\textsuperscript{o}2023-150 du 12 avril 2023)}{fig:organigramme}

\section{Travaux effectu\'es}

Au sein de la Cellule de la Statistique, nous avons particip\'e aux activit\'es
courantes de collecte, de mise \`a jour et de pr\'earchivage des documents et
donn\'ees statistiques de la direction, ce qui nous a permis de d\'ecouvrir
concr\`etement les m\'ecanismes de suivi des activit\'es de recherche et
d'innovation au niveau national.

L'essentiel de notre temps de stage a toutefois \'et\'e consacr\'e \`a la
conception et au d\'eveloppement de notre projet de programmation,
\textbf{CyberGuard B\'enin}, r\'ealis\'e dans le cadre de notre m\'emoire de fin de
formation. Ce travail s'est d\'eroul\'e en plusieurs \'etapes~:

\begin{itemize}[nosep]
  \item clarification du besoin et des objectifs du projet~;
  \item choix de la stack technique (Laravel pour l'API backend, React pour le
    frontend, MySQL comme syst\`eme de gestion de base de donn\'ees)~;
  \item mod\'elisation UML du syst\`eme~: identification des acteurs et des cas
    d'utilisation, diagrammes de classes, d'objets et de s\'equence~;
  \item conception de la base de donn\'ees et d\'efinition du dictionnaire des
    donn\'ees~;
  \item d\'efinition puis red\'efinition compl\`ete de l'identit\'e visuelle de la
    plateforme, suivie de la r\'ealisation de maquettes haute-fid\'elit\'e~;
  \item impl\'ementation du moteur d'analyse de s\'ecurit\'e (sept contr\^oles non
    intrusifs~: HTTPS, certificat SSL, en-t\^etes de s\'ecurit\'e, cookies,
    vuln\'erabilit\'es, fichiers sensibles expos\'es, s\'ecurit\'e email SPF/DMARC),
    du calcul du score et des recommandations~;
  \item impl\'ementation de l'authentification (par email/mot de passe et via
    Google), du tableau de bord, de l'historique des analyses et des
    notifications in-app~;
  \item impl\'ementation du back-office d'administration (gestion des comptes,
    mod\'eration des analyses, param\'etrage du bar\`eme de scoring)~;
  \item tests fonctionnels et corrections successives, jusqu'\`a l'obtention
    d'une plateforme stable et utilisable.
\end{itemize}

\section{Apport du stage sur le plan professionnel}

Ce stage nous a permis de mettre en pratique, sur un projet r\'eel et de bout en
bout, l'ensemble du cycle de vie d'un projet informatique~: de la clarification
du besoin jusqu'\`a une plateforme fonctionnelle, en passant par la
mod\'elisation UML, la conception de base de donn\'ees et le d\'eveloppement
technique. Il a \'egalement renforc\'e nos comp\'etences en d\'eveloppement web
full-stack (Laravel, React), en s\'ecurit\'e applicative (authentification,
protection contre le sondage d'adresses r\'eseau internes, validation des
entr\'ees), ainsi qu'en gestion it\'erative d'un projet~: priorisation des
fonctionnalit\'es, correction de r\'egressions, et adaptation continue de la
plateforme aux retours obtenus au fil du d\'eveloppement. Le passage par la
Cellule de la Statistique nous a par ailleurs sensibilis\'es \`a l'importance
d'une donn\'ee fiable et \`a jour pour le pilotage de politiques publiques -- un
constat qui fait \'echo \`a la probl\'ematique m\^eme de CyberGuard B\'enin~: donner
\`a des acteurs non-techniciens une information fiable et compr\'ehensible sur
laquelle agir.

\section{Difficult\'es rencontr\'ees}

Le d\'eveloppement de CyberGuard B\'enin n'a pas \'et\'e exempt de difficult\'es
techniques. Nous en retenons notamment quatre~:

\begin{itemize}[nosep]
  \item \textbf{Authentification via Google et gestion de session}~:
    l'int\'egration de la connexion via Google (Laravel Socialite) s'est
    heurt\'ee \`a un comportement du middleware de session de Laravel Sanctum,
    qui ne d\'emarre une session que lorsque la requ\^ete est reconnue comme
    provenant du frontend (via l'en-t\^ete \emph{Referer}) -- ce qui n'est pas le
    cas d'un retour depuis Google. Il a fallu forcer explicitement le d\'emarrage
    de session sur les routes concern\'ees.
  \item \textbf{S\'ecurit\'e du moteur d'analyse}~: le moteur devant accepter une
    URL saisie librement par l'utilisateur, une attention particuli\`ere a \'et\'e
    port\'ee \`a emp\^echer qu'il ne soit d\'etourn\'e pour sonder le r\'eseau interne
    du serveur (SSRF)~; la r\'esolution DNS de l'h\^ote est donc v\'erifi\'ee \`a deux
    reprises (\`a la soumission puis juste avant l'analyse effective).
  \item \textbf{Fiabilit\'e du traitement en arri\`ere-plan}~: les analyses \'etant
    trait\'ees de mani\`ere asynchrone via une file d'attente, un oubli de
    d\'emarrage du worker de traitement laissait certaines analyses bloqu\'ees
    ind\'efiniment au statut \og~en cours~\fg{} -- un point de vigilance
    r\'ecurrent en environnement de d\'eveloppement local.
  \item \textbf{Accessibilit\'e des r\'esultats pour un public non-technicien}~:
    conform\'ement \`a l'objectif du projet, les premiers constats techniques
    g\'en\'er\'es (en-t\^etes HTTP manquants, versions de protocole obsol\`etes...)
    se sont r\'ev\'el\'es trop jargonneux pour l'utilisateur cible. L'ensemble des
    descriptions et recommandations a d\^u \^etre reformul\'e en langage courant,
    compr\'ehensible sans connaissance pr\'ealable en s\'ecurit\'e informatique.
\end{itemize}

Ces difficult\'es, une fois identifi\'ees, ont chacune \'et\'e corrig\'ees et
v\'erifi\'ees avant la poursuite du d\'eveloppement.

\chapter{Projet de programmation}
\label{chap:programmation}

\section{Pr\'esentation du projet de programmation}

La v\'erification de la s\'ecurit\'e d'un site web constitue un probl\`eme majeur pour
les entreprises, \'ecoles et administrations b\'eninoises, particuli\`erement dans un
contexte o\`u l'expertise en cybers\'ecurit\'e reste rare et co\^uteuse. Les m\'ethodes
disponibles aujourd'hui -- audits professionnels payants, outils techniques
disperss\'es et en anglais, ou absence totale de v\'erification -- ne r\'epondent pas
au besoin d'un propri\'etaire de site qui veut simplement savoir si son site est
s\^ur et quoi corriger. \textbf{CyberGuard B\'enin} est une plateforme web con\c{c}ue
pour rendre ce diagnostic accessible, automatique et compr\'ehensible.

\section{Pr\'esentation de la solution}

Pour r\'epondre \`a la probl\'ematique de d\'emocratisation de l'audit de s\'ecurit\'e web,
nous proposons le d\'eveloppement d'une plateforme compl\`ete utilisant les
technologies Laravel et React. L'utilisateur cr\'ee un compte, saisit l'URL de son
site, et la plateforme lance une s\'erie de v\'erifications non intrusives (HTTPS,
certificat SSL, en-t\^etes de s\'ecurit\'e, exposition d'informations serveur,
cookies). Le syst\`eme traduit ces r\'esultats techniques en un score sur 100 et un
niveau de risque (Bon / Moyen / Faible), accompagn\'e de recommandations concr\`etes
et compr\'ehensibles. Un rapport est g\'en\'er\'e, et l'utilisateur peut suivre
l'\'evolution de la s\'ecurit\'e de son site dans le temps via un tableau de bord et
un historique.

\section{Analyse et mod\'elisation}

\subsection{Choix technique}

La mod\'elisation de ce projet s'appuie sur les diagrammes UML standard (cas
d'utilisation, classes, objets, s\'equence), qui permettent de visualiser les
acteurs, les interactions et la structure statique du syst\`eme avant son
impl\'ementation. Le tableau~\ref{tab:stack} r\'esume les choix technologiques
retenus.

\begin{table}[H]
  \centering
  \caption{Stack technique retenue}
  \label{tab:stack}
  \begin{tabular}{ll}
    \toprule
    \textbf{Couche} & \textbf{Choix} \\
    \midrule
    Frontend & React (Vite) \\
    Backend & Laravel \\
    Base de donn\'ees & MySQL (MariaDB via XAMPP) \\
    Authentification & Laravel Sanctum \\
    G\'en\'eration PDF & DomPDF \\
    \bottomrule
  \end{tabular}
\end{table}

\subsection{\'Etude de l'existant}

Des outils comme \textbf{Mozilla Observatory} et \textbf{SSL Labs} existent d\'ej\`a
pour analyser la s\'ecurit\'e d'un site web. Ils sont reconnus et techniquement
solides, en s'appuyant sur des v\'erifications approfondies des en-t\^etes de
s\'ecurit\'e et des certificats SSL/TLS.

\subsection{Insuffisances}

Ces outils pr\'esentent toutefois plusieurs limites dans le contexte vis\'e par
notre projet. D'abord, ils sont exclusivement en anglais, ce qui constitue un
frein r\'eel pour une partie du public b\'eninois. Ensuite, leurs r\'esultats restent
tr\`es techniques~: un propri\'etaire de site non-technicien obtient une note ou une
liste d'en-t\^etes manquants, sans traduction claire de ce que cela signifie
concr\`etement ni de la marche \`a suivre pour corriger le probl\`eme. Enfin, ces
outils sont dispers\'es~: v\'erifier HTTPS, le certificat SSL et les en-t\^etes de
s\'ecurit\'e n\'ecessite souvent de passer par plusieurs services diff\'erents, sans
suivi centralis\'e ni historique des analyses pass\'ees. CyberGuard B\'enin se
positionne comme une alternative qui centralise ces v\'erifications, restitue un
score unique et des recommandations en fran\c{c}ais compr\'ehensibles par un public
non-technicien, tout en conservant un historique consultable dans le temps.

\subsection{Repr\'esentation de la mod\'elisation}

\subsubsection{Diagramme des cas d'utilisation}

\paragraph{Identification des acteurs}

Trois acteurs interagissent avec le syst\`eme (tableau~\ref{tab:acteurs}).

\begin{table}[H]
  \centering
  \caption{Identification des acteurs}
  \label{tab:acteurs}
  \begin{tabular}{lp{9.5cm}}
    \toprule
    \textbf{Acteur} & \textbf{Description} \\
    \midrule
    Utilisateur & Personne disposant d'un compte sur la plateforme, propri\'etaire ou responsable d'un site web qu'elle souhaite analyser. \\
    Administrateur & Personne charg\'ee de la gestion de la plateforme (comptes, statistiques, param\`etres de scoring). H\'erite de tous les droits de l'Utilisateur. \\
    Syst\`eme & Acteur non-humain~: partie automatis\'ee de l'application qui ex\'ecute les v\'erifications techniques et produit les r\'esultats. \\
    \bottomrule
  \end{tabular}
\end{table}

\paragraph{Identification des cas d'utilisation}

Pour chaque acteur, les cas d'utilisation suivants ont \'et\'e identifi\'es~:

\textbf{Utilisateur}
\begin{itemize}[nosep]
  \item G\'erer son compte (inscription, modification, r\'einitialisation du mot de passe)
  \item S'authentifier (connexion par email/mot de passe, ou via un compte Google d\'ej\`a existant -- Google n'est jamais un moyen d'inscription)
  \item Lancer une analyse
  \item Consulter les rapports (historique, d\'etail, t\'el\'echargement PDF)
  \item Consulter le tableau de bord
  \item \^Etre notifi\'e de la fin d'une analyse (notification in-app lorsqu'une analyse lanc\'ee en arri\`ere-plan se termine)
\end{itemize}
\textbf{Administrateur}
\begin{itemize}[nosep]
  \item G\'erer les comptes utilisateurs (activer/d\'esactiver, promouvoir/r\'etrograder le r\^ole, supprimer)
  \item Mod\'erer les analyses (consulter et supprimer l'analyse de n'importe quel utilisateur)
  \item Administrer la plateforme (statistiques globales, param\`etres de scoring)
\end{itemize}
\textbf{Syst\`eme}
\begin{itemize}[nosep]
  \item Traiter l'analyse (v\'erifications techniques, calcul du score, g\'en\'eration du rapport, notification de l'utilisateur)
\end{itemize}

\paragraph{Tableau r\'ecapitulatif des cas d'utilisation}

\begin{longtable}{clc}
  \caption{Tableau r\'ecapitulatif des cas d'utilisation}
  \label{tab:recap-cu} \\
  \toprule
  \textbf{\#} & \textbf{Cas d'utilisation} & \textbf{Acteur} \\
  \midrule
  \endfirsthead
  \toprule
  \textbf{\#} & \textbf{Cas d'utilisation} & \textbf{Acteur} \\
  \midrule
  \endhead
  1 & G\'erer son compte & Utilisateur \\
  2 & S'authentifier & Utilisateur \\
  3 & Lancer une analyse & Utilisateur \\
  4 & Consulter les rapports & Utilisateur \\
  5 & Consulter le tableau de bord & Utilisateur \\
  6 & \^Etre notifi\'e de la fin d'une analyse & Utilisateur \\
  7 & G\'erer les comptes utilisateurs & Administrateur \\
  8 & Mod\'erer les analyses & Administrateur \\
  9 & Administrer la plateforme & Administrateur \\
  10 & Traiter l'analyse & Syst\`eme \\
  \bottomrule
\end{longtable}

\paragraph{Diagramme proprement dit}

Le diagramme ci-dessous formalise les acteurs et cas d'utilisation identifi\'es,
avec une relation de g\'en\'eralisation (Administrateur h\'erite d'Utilisateur) et
une relation d'inclusion (\og~Lancer une analyse~\fg{} inclut \og~Traiter
l'analyse~\fg).

\TODOFigure{diagramme-cas-utilisation.jpeg}{Diagramme de cas d'utilisation}{fig:cas-utilisation}

\paragraph{Description textuelle de cinq cas d'utilisation}

\subparagraph{Cas d'utilisation~: G\'erer les comptes utilisateurs}
\begin{table}[H]
  \centering
  \small
  \begin{tabular}{p{3.2cm}p{9.5cm}}
    \toprule
    Acteur principal & Administrateur \\
    Objectif & Permettre \`a un administrateur de g\'erer les comptes des utilisateurs de la plateforme (statut, r\^ole). \\
    Pr\'econdition & L'administrateur est authentifi\'e. \\
    Sc\'enario nominal &
      1. Acc\`es \`a la liste des comptes utilisateurs.
      2. Recherche par nom/email, filtre par statut.
      3. Activation ou d\'esactivation d'un compte.
      4. Promotion ou r\'etrogradation du r\^ole (utilisateur $\leftrightarrow$ administrateur).
      5. Suppression d'un compte si n\'ecessaire. \\
    Postcondition & L'\'etat du compte cibl\'e (statut, r\^ole, ou existence) est mis \`a jour en cons\'equence. \\
    Exception & Tentative d'un administrateur de d\'esactiver, r\'etrograder ou supprimer son propre compte $\rightarrow$ refus\'ee (garde-fou anti-auto-modification). Tentative d'acc\`es par un utilisateur non administrateur $\rightarrow$ acc\`es refus\'e. \\
    \bottomrule
  \end{tabular}
\end{table}

\subparagraph{Cas d'utilisation~: Lancer une analyse}
\begin{table}[H]
  \centering
  \small
  \begin{tabular}{p{3.2cm}p{9.5cm}}
    \toprule
    Acteur principal & Utilisateur \\
    Objectif & Permettre \`a l'utilisateur de demander l'analyse de s\'ecurit\'e d'un site web. \\
    Pr\'econdition & L'utilisateur est authentifi\'e. \\
    Sc\'enario nominal &
      1. L'utilisateur saisit l'URL du site \`a analyser.
      2. Le syst\`eme valide le format de l'URL.
      3. Le syst\`eme d\'eclenche le cas d'utilisation \og~Traiter l'analyse~\fg.
      4. Une confirmation de lancement est affich\'ee \`a l'utilisateur. \\
    Postcondition & Une analyse est en cours de traitement pour le site indiqu\'e. \\
    Exception & URL invalide ou site inaccessible $\rightarrow$ message d'erreur affich\'e, aucune analyse n'est lanc\'ee. \\
    \bottomrule
  \end{tabular}
\end{table}

\subparagraph{Cas d'utilisation~: Traiter l'analyse}
\begin{table}[H]
  \centering
  \small
  \begin{tabular}{p{3.2cm}p{9.5cm}}
    \toprule
    Acteur principal & Syst\`eme \\
    Objectif & Ex\'ecuter automatiquement les v\'erifications de s\'ecurit\'e et produire un r\'esultat exploitable. \\
    Pr\'econdition & Le cas \og~Lancer une analyse~\fg{} a \'et\'e d\'eclench\'e avec une URL valide. \\
    Sc\'enario nominal &
      1. V\'erification que l'h\^ote ne pointe pas vers une ressource r\'eseau priv\'ee ou interne.
      2. V\'erification de la pr\'esence du HTTPS.
      3. Contr\^ole de la validit\'e du certificat SSL.
      4. Analyse des en-t\^etes de s\'ecurit\'e HTTP.
      5. V\'erification de la configuration des cookies (Secure, HttpOnly, SameSite).
      6. Recherche de vuln\'erabilit\'es simples et non intrusives.
      7. Recherche de fichiers sensibles expos\'es publiquement (.env, .git...).
      8. V\'erification de la protection du domaine contre l'usurpation d'email (SPF/DMARC).
      9. Calcul du score sur 100 et du niveau de risque.
      10. G\'en\'eration des recommandations, en langage compr\'ehensible par un non-technicien.
      11. G\'en\'eration du rapport PDF.
      12. Enregistrement du r\'esultat en base de donn\'ees.
      13. Notification de l'utilisateur si l'analyse \'etait suivie en arri\`ere-plan (\og~\^Etre notifi\'e de la fin d'une analyse~\fg). \\
    Postcondition & Un rapport d'analyse (score, risque, recommandations, PDF) est disponible et associ\'e \`a l'utilisateur. \\
    Exception & Adresse non autoris\'ee, site inaccessible ou timeout $\rightarrow$ l'analyse est marqu\'ee en \'echec avec un motif explicite, l'utilisateur est notifi\'e. \\
    \bottomrule
  \end{tabular}
\end{table}

\subparagraph{Cas d'utilisation~: Consulter les rapports}
\begin{table}[H]
  \centering
  \small
  \begin{tabular}{p{3.2cm}p{9.5cm}}
    \toprule
    Acteur principal & Utilisateur \\
    Objectif & Permettre \`a l'utilisateur de revoir ses analyses pass\'ees et d'en r\'ecup\'erer le d\'etail. \\
    Pr\'econdition & L'utilisateur est authentifi\'e et poss\`ede au moins une analyse termin\'ee. \\
    Sc\'enario nominal &
      1. L'utilisateur acc\`ede \`a son historique d'analyses.
      2. Le syst\`eme affiche la liste des analyses (date, site, score, niveau de risque).
      3. L'utilisateur s\'electionne une analyse.
      4. Le syst\`eme affiche le d\'etail du rapport.
      5. L'utilisateur peut t\'el\'echarger le rapport au format PDF. \\
    Postcondition & L'utilisateur a consult\'e ou t\'el\'echarg\'e le rapport souhait\'e. \\
    Exception & Aucune analyse enregistr\'ee $\rightarrow$ message indiquant qu'aucun historique n'est disponible. \\
    \bottomrule
  \end{tabular}
\end{table}

\subparagraph{Cas d'utilisation~: Mod\'erer les analyses}
\begin{table}[H]
  \centering
  \small
  \begin{tabular}{p{3.2cm}p{9.5cm}}
    \toprule
    Acteur principal & Administrateur \\
    Objectif & Permettre \`a un administrateur de superviser l'ensemble des analyses de la plateforme, tous utilisateurs confondus, et d'en supprimer une si n\'ecessaire. \\
    Pr\'econdition & L'administrateur est authentifi\'e. \\
    Sc\'enario nominal &
      1. Acc\`es \`a la liste globale des analyses.
      2. Affichage de toutes les analyses avec leur propri\'etaire, ind\'ependamment de l'utilisateur connect\'e.
      3. Recherche par URL ou par propri\'etaire, filtre par niveau de risque.
      4. Consultation du d\'etail d'une analyse.
      5. Suppression de l'analyse (et de son rapport PDF associ\'e) si n\'ecessaire. \\
    Postcondition & La liste refl\`ete l'\'etat courant des analyses~; une analyse supprim\'ee n'est plus consultable. \\
    Exception & Tentative d'acc\`es par un utilisateur non administrateur $\rightarrow$ acc\`es refus\'e. \\
    \bottomrule
  \end{tabular}
\end{table}

\subsubsection{Diagramme des classes}

Un diagramme de classes est un diagramme UML qui repr\'esente la structure
statique d'un syst\`eme logiciel en mod\'elisant les classes avec leurs attributs,
m\'ethodes et les relations qui les lient. Chaque classe est visualis\'ee sous forme
d'un rectangle divis\'e en trois sections (nom, attributs, m\'ethodes). Les relations
entre classes (association, h\'eritage, agr\'egation, d\'ependance) sont
repr\'esent\'ees par diff\'erents types de fl\`eches et de lignes.

\paragraph{Identification des classes}
\begin{itemize}
  \item Utilisateur
  \item Administrateur
  \item Analyse
  \item Controle
  \item ParametreScore
\end{itemize}

\paragraph{R\`egles de gestion}
\textbf{Utilisateur}
\begin{itemize}[nosep]
  \item Un Utilisateur peut lancer plusieurs Analyses.
  \item Un Administrateur h\'erite de tous les droits d'un Utilisateur (g\'en\'eralisation).
\end{itemize}
\textbf{Analyse}
\begin{itemize}[nosep]
  \item Une Analyse est lanc\'ee par un seul Utilisateur.
  \item Une Analyse contient plusieurs Controles (r\'esultats de v\'erification).
  \item Une Analyse a un statut \'evolutif (en\_cours, terminee, echec).
  \item Un Administrateur peut consulter et supprimer n'importe quelle Analyse, sans restriction de propri\'etaire (mod\'eration).
\end{itemize}
\textbf{Controle}
\begin{itemize}[nosep]
  \item Un Controle est toujours li\'e \`a une seule Analyse.
  \item Un Controle a un statut (conforme / non\_conforme) et peut porter une recommandation.
\end{itemize}
\textbf{ParametreScore}
\begin{itemize}[nosep]
  \item Il n'existe qu'une seule instance de ParametreScore (configuration globale, singleton).
  \item Un Administrateur peut configurer les param\`etres de score (\og~configure~\fg).
\end{itemize}

\paragraph{Dictionnaire des donn\'ees}

Le dictionnaire des donn\'ees ci-dessous recense, pour chaque table de la base de
donn\'ees, l'ensemble des champs, leur type SQL, leur taille et leur r\^ole (cf.
\texttt{schema.sql}).

\begin{table}[H]
  \centering
  \caption{Dictionnaire des donn\'ees -- table \texttt{utilisateurs}}
  \begin{tabular}{lllp{6cm}}
    \toprule
    \textbf{Champ} & \textbf{Type} & \textbf{Taille} & \textbf{Description} \\
    \midrule
    id & BIGINT UNSIGNED & -- & Identifiant unique de l'utilisateur \\
    nom & VARCHAR & 150 & Nom complet de l'utilisateur \\
    email & VARCHAR & 190 & Adresse email (identifiant de connexion) \\
    mot\_de\_passe & VARCHAR & 255 & Mot de passe hach\'e \\
    role & ENUM & -- & \texttt{utilisateur} / \texttt{administrateur} \\
    statut & ENUM & -- & \texttt{actif} / \texttt{desactive} \\
    date\_inscription & TIMESTAMP & -- & Date de cr\'eation du compte \\
    \bottomrule
  \end{tabular}
\end{table}

\begin{table}[H]
  \centering
  \caption{Dictionnaire des donn\'ees -- table \texttt{analyses}}
  \begin{tabular}{lllp{6cm}}
    \toprule
    \textbf{Champ} & \textbf{Type} & \textbf{Taille} & \textbf{Description} \\
    \midrule
    id & BIGINT UNSIGNED & -- & Identifiant unique de l'analyse \\
    utilisateur\_id & BIGINT UNSIGNED & -- & R\'ef\'erence vers l'utilisateur propri\'etaire \\
    url & VARCHAR & 255 & URL du site analys\'e \\
    date\_analyse & TIMESTAMP & -- & Date de lancement de l'analyse \\
    score & TINYINT UNSIGNED & -- & Score obtenu sur 100 \\
    niveau\_risque & ENUM & -- & \texttt{bon} / \texttt{moyen} / \texttt{faible} \\
    statut & ENUM & -- & \texttt{en\_cours} / \texttt{terminee} / \texttt{echec} \\
    chemin\_rapport\_pdf & VARCHAR & 255 & Chemin du rapport PDF g\'en\'er\'e \\
    \bottomrule
  \end{tabular}
\end{table}

\begin{table}[H]
  \centering
  \caption{Dictionnaire des donn\'ees -- table \texttt{controles}}
  \begin{tabular}{lllp{6cm}}
    \toprule
    \textbf{Champ} & \textbf{Type} & \textbf{Taille} & \textbf{Description} \\
    \midrule
    id & BIGINT UNSIGNED & -- & Identifiant unique du contr\^ole \\
    analyse\_id & BIGINT UNSIGNED & -- & R\'ef\'erence vers l'analyse concern\'ee \\
    type & VARCHAR & 50 & https, certificat\_ssl, en\_tetes\_securite, cookies, vulnerabilites, fichiers\_exposes, email\_securise \\
    statut & ENUM & -- & \texttt{conforme} / \texttt{non\_conforme} \\
    description & TEXT & -- & D\'etail du constat \\
    recommandation & TEXT & -- & Recommandation de correction (si non conforme) \\
    \bottomrule
  \end{tabular}
\end{table}

\begin{table}[H]
  \centering
  \caption{Dictionnaire des donn\'ees -- table \texttt{parametres\_score}}
  \begin{tabular}{lllp{6cm}}
    \toprule
    \textbf{Champ} & \textbf{Type} & \textbf{Taille} & \textbf{Description} \\
    \midrule
    id & TINYINT UNSIGNED & -- & Identifiant fixe (toujours 1, ligne unique) \\
    poids\_https & TINYINT UNSIGNED & -- & Pond\'eration du crit\`ere HTTPS \\
    poids\_certificat\_ssl & TINYINT UNSIGNED & -- & Pond\'eration du crit\`ere certificat SSL \\
    poids\_headers & TINYINT UNSIGNED & -- & Pond\'eration du crit\`ere en-t\^etes \\
    poids\_cookies & TINYINT UNSIGNED & -- & Pond\'eration du crit\`ere cookies \\
    poids\_vulnerabilites & TINYINT UNSIGNED & -- & Pond\'eration du crit\`ere vuln\'erabilit\'es \\
    poids\_fichiers\_exposes & TINYINT UNSIGNED & -- & Pond\'eration du crit\`ere fichiers expos\'es \\
    poids\_email\_securise & TINYINT UNSIGNED & -- & Pond\'eration du crit\`ere s\'ecurit\'e email (SPF/DMARC) \\
    seuil\_bon & TINYINT UNSIGNED & -- & Seuil de score pour le niveau \og~bon~\fg \\
    seuil\_moyen & TINYINT UNSIGNED & -- & Seuil de score pour le niveau \og~moyen~\fg \\
    \bottomrule
  \end{tabular}
\end{table}

\paragraph{Diagramme de classes}

\TODOFigure{diagramme-classes.jpeg}{Diagramme de classes}{fig:classes}

\subsubsection{Diagramme d'objet}

Le diagramme d'objets ci-dessous illustre une instanciation concr\`ete du
diagramme de classes~: un utilisateur ayant lanc\'e une analyse termin\'ee avec
trois contr\^oles, et le bar\`eme de score en vigueur, configur\'e par un
administrateur.

\TODOFigure{diagramme-objets.jpeg}{Diagramme d'objets}{fig:objets}

\subsubsection{Diagrammes de s\'equence}

Un diagramme de s\'equence par cas d'utilisation d\'ecrit ci-dessus, entre
l'Utilisateur, le Syst\`eme et la base de donn\'ees.

\paragraph{G\'erer les comptes utilisateurs}

Liste des comptes avec recherche et filtre par statut, puis modification
optionnelle (statut, r\^ole ou suppression), prot\'eg\'ee par un garde-fou
emp\^echant un administrateur de modifier son propre compte.

\TODOFigure{diagramme-sequence-gerer-comptes.jpeg}{Diagramme de s\'equence -- G\'erer les comptes utilisateurs}{fig:sequence-gerer-comptes}

\paragraph{Lancer une analyse}

Saisie de l'URL, validation du format, cr\'eation de l'analyse (statut
\texttt{en\_cours}) en base et confirmation \`a l'utilisateur, avec le chemin
d'exception \og~URL invalide~\fg. Ce cas d\'eclenche automatiquement \og~Traiter
l'analyse~\fg{} (relation \og~include~\fg).

\TODOFigure{diagramme-sequence-lancer-analyse.jpeg}{Diagramme de s\'equence -- Lancer une analyse}{fig:sequence-lancer}

\paragraph{Traiter l'analyse}

Ex\'ecution des sept contr\^oles (HTTPS, certificat SSL, en-t\^etes de
s\'ecurit\'e, cookies, vuln\'erabilit\'es, fichiers expos\'es, s\'ecurit\'e email
SPF/DMARC), calcul du score et du niveau de risque, g\'en\'eration des
recommandations et du rapport, puis enregistrement en base, avec le chemin
d'exception \og~adresse non autoris\'ee, site inaccessible ou timeout~\fg. Le
traitement est d\'el\'egu\'e \`a une file d'attente~: une fois termin\'e,
l'utilisateur est notifi\'e dans l'application s'il suivait cette analyse
(\og~\^Etre notifi\'e de la fin d'une analyse~\fg).

\TODOFigure{diagramme-sequence-traiter-analyse.jpeg}{Diagramme de s\'equence -- Traiter l'analyse}{fig:sequence-traiter}

\paragraph{Consulter les rapports}

Acc\`es \`a l'historique des analyses, s\'election d'une analyse pour en consulter
le d\'etail, et t\'el\'echargement optionnel du rapport PDF, avec le chemin
d'exception \og~aucune analyse enregistr\'ee~\fg.

\paragraph{Mod\'erer les analyses}

Vue globale r\'eserv\'ee \`a l'Administrateur~: m\^emes principes de recherche et
de filtrage que \og~Consulter les rapports~\fg, mais sans restriction au
propri\'etaire, avec une action de suppression suppl\'ementaire. Acc\`es prot\'eg\'e
par le r\^ole administrateur~: une tentative par un utilisateur normal est
refus\'ee (403).

\section{Outils utilis\'es}

\subsection{Mat\'eriels utilis\'es}

Le d\'eveloppement a \'et\'e r\'ealis\'e sur un ordinateur portable \textbf{HP ProBook
450 G7}, dot\'e d'un processeur \textbf{Intel Core i7-10510U} cadenc\'e \`a
1,80~GHz (jusqu'\`a 2,30~GHz en turbo), de \textbf{16~Go de m\'emoire RAM}, d'un
disque de stockage de 238~Go et d'une carte graphique int\'egr\'ee Intel UHD
Graphics, sous \textbf{Windows 11 Pro}.

\subsection{Frontend}
\begin{itemize}
  \item \textbf{React (Vite)}~: biblioth\`eque JavaScript utilis\'ee pour construire
    une interface utilisateur dynamique en Single Page Application, avec
    \texttt{react-router-dom} pour la navigation et \texttt{axios} pour la
    communication avec l'API.
\end{itemize}

\subsection{Backend (API et logique m\'etier)}
\begin{itemize}
  \item \textbf{Laravel}~: framework PHP servant de base \`a l'API RESTful,
    responsable de la logique m\'etier (moteur de v\'erification de s\'ecurit\'e), de
    l'interaction avec la base de donn\'ees via l'ORM Eloquent, et de la gestion
    de l'authentification via \textbf{Laravel Sanctum}.
\end{itemize}

\subsection{Environnement de d\'eveloppement local}
\begin{itemize}
  \item \textbf{XAMPP}~: environnement logiciel int\'egr\'e fournissant Apache, PHP
    et MySQL (MariaDB) pour le d\'eveloppement local de l'application Laravel.
\end{itemize}

\subsection{Conception d'interface}
\begin{itemize}
  \item \textbf{Stitch} (Google)~: outil de g\'en\'eration de maquettes haute-fid\'elit\'e \`a
    partir de prompts textuels, utilis\'e pour produire les \'ecrans pr\'esent\'es en
    section~\ref{sec:maquettes} \`a partir de l'identit\'e visuelle d\'efinie pour le
    projet (palette de couleurs, typographie, style des composants).
\end{itemize}

\subsection{Outils de test}

Deux niveaux de test compl\'ementaires ont \'et\'e utilis\'es tout au long du
d\'eveloppement.

\paragraph{Tests automatis\'es (backend)} L'API est couverte par une suite de
tests \textbf{PHPUnit}, ex\'ecut\'ee via la commande \texttt{php artisan test}.
Chaque test s'appuie sur le trait \texttt{RefreshDatabase} pour repartir d'une
base de donn\'ees propre, et la biblioth\`eque \textbf{Mockery} est utilis\'ee
pour simuler les r\'eponses du service tiers Google (fa\c{c}ade Socialite) sans
d\'ependre d'une v\'eritable connexion \`a Google lors des tests. Ces tests
couvrent notamment~: l'authentification (classique et via Google), les
autorisations d'acc\`es (propri\'etaire d'une ressource, r\^ole administrateur,
garde-fous anti-auto-modification), la validation des URL soumises \`a
analyse (rejet des adresses r\'eseau priv\'ees ou invalides), et les r\`egles de
gestion des param\`etres de scoring.

\paragraph{Tests manuels de bout en bout (frontend)} L'interface a \'et\'e
v\'erifi\'ee de mani\`ere r\'ep\'et\'ee \`a l'aide de \textbf{Playwright}, un outil
d'automatisation de navigateur~: simulation d'un parcours utilisateur complet
(connexion, lancement d'une analyse, consultation du r\'esultat), v\'erification
de l'absence de d\'ebordement horizontal sur mobile \`a plusieurs largeurs
d'\'ecran, et contr\^ole du bon fonctionnement des fonctionnalit\'es
d'administration. Cette approche a permis de d\'etecter plusieurs anomalies
r\'eelles avant qu'elles n'affectent l'utilisateur final (par exemple, une
session non d\'emarr\'ee sur le chemin de retour de l'authentification Google, ou
une condition de course dans le suivi des notifications en arri\`ere-plan).

\subsection{D\'eploiement (mise en production)}

\`A ce stade du projet, la plateforme n'est pas encore d\'eploy\'ee en production~:
le d\'eveloppement et les tests se font en environnement local via XAMPP
(Apache, MySQL) pour le backend et le serveur de d\'eveloppement Vite pour le
frontend. Le choix d'un h\'ebergeur et la mise en production font partie des
perspectives d'\'evolution du projet (voir conclusion).

\section{Maquettes}
\label{sec:maquettes}

Sur la base de l'identit\'e visuelle d\'efinie pour CyberGuard B\'enin --
\og~Le Sceau~\fg, une identit\'e de confiance institutionnelle~: papier chaud
(\#F5F1E6), bleu marine (\#16304A), accent or (\#B8863A), typographie
Georgia pour les titres, composants aux coins nets et bordures fines plut\^ot
que des ombres port\'ees -- des maquettes haute-fid\'elit\'e ont \'et\'e
r\'ealis\'ees avec \textbf{Stitch} pour chacun des \'ecrans identifi\'es \`a
partir des cas d'utilisation, avant l'impl\'ementation compl\`ete du frontend.
L'identit\'e visuelle a \'et\'e enti\`erement red\'efinie en cours de projet~; les
maquettes ci-dessous refl\`etent la version actuelle.

\TODOFigure{maquette-accueil.png}{Accueil, avec la modale Connexion/Inscription}{fig:maq-accueil}
\TODOFigure{maquette-mot-de-passe-oublie.png}{Mot de passe oubli\'e}{fig:maq-mdp-oublie}
\TODOFigure{maquette-mon-profil.png}{Mon profil}{fig:maq-profil}
\TODOFigure{maquette-tableau-de-bord-mobile.png}{Tableau de bord (variante mobile)}{fig:maq-dashboard-mobile}
\TODOFigure{maquette-tableau-de-bord-desktop.png}{Tableau de bord (variante desktop)}{fig:maq-dashboard-desktop}
\TODOFigure{maquette-nouvelle-analyse.png}{Nouvelle analyse (analyse en cours)}{fig:maq-nouvelle-analyse}
\TODOFigure{maquette-resultat-analyse.png}{R\'esultat d'analyse}{fig:maq-resultat}
\TODOFigure{maquette-historique.png}{Historique des analyses}{fig:maq-historique}
\TODOFigure{maquette-detail-rapport.png}{D\'etail d'un rapport, avec t\'el\'echargement PDF}{fig:maq-detail-rapport}
\TODOFigure{maquette-admin-utilisateurs.png}{Administration -- gestion des utilisateurs}{fig:maq-admin-utilisateurs}
\TODOFigure{maquette-admin-analyses.png}{Administration -- mod\'eration des analyses}{fig:maq-admin-analyses}
\TODOFigure{maquette-admin-parametres.png}{Administration -- statistiques et param\`etres}{fig:maq-admin-parametres}

\begin{table}[H]
  \centering
  \caption{Correspondance maquette / cas d'utilisation}
  \begin{tabular}{p{6cm}p{6.5cm}}
    \toprule
    \textbf{Maquette} & \textbf{Cas d'utilisation associ\'e} \\
    \midrule
    Accueil & S'authentifier (email/mot de passe ou Google), G\'erer son compte (inscription) \\
    Mot de passe oubli\'e & G\'erer son compte (r\'einitialisation) \\
    Mon profil & G\'erer son compte (modification) \\
    Tableau de bord & Consulter le tableau de bord \\
    Nouvelle analyse & Lancer une analyse \\
    R\'esultat d'analyse & Traiter l'analyse (r\'esultat), \^Etre notifi\'e de la fin d'une analyse \\
    Historique & Consulter les rapports (liste) \\
    D\'etail d'un rapport & Consulter les rapports (d\'etail + PDF) \\
    Administration -- utilisateurs & G\'erer les comptes utilisateurs \\
    Administration -- analyses & Mod\'erer les analyses \\
    Administration -- param\`etres & Administrer la plateforme \\
    \bottomrule
  \end{tabular}
\end{table}

\chapter{R\'esultats obtenus}
\label{chap:resultats}

Ce chapitre pr\'esente les r\'esultats techniques obtenus lors du d\'eveloppement de
notre application~: l'architecture mise en place, les mesures de s\'ecurit\'e
appliqu\'ees, ainsi que certaines interfaces illustrant le fonctionnement de la
plateforme.

\section{R\'esultats}

\subsection{Pr\'esentation de quelques interfaces de l'application}

% TODO : remplacer chaque \TODOFigure par une vraie capture d'\'ecran une fois
% le frontend d\'evelopp\'e (page d'accueil, tableau de bord, nouvelle analyse,
% r\'esultat d'analyse, historique, administration...).

\subsubsection{Page d'accueil}
\textit{[\`A compl\'eter]}
\TODOFigure{page-accueil.png}{Page d'accueil}{fig:accueil}

\subsubsection{Tableau de bord}
\textit{[\`A compl\'eter]}
\TODOFigure{tableau-de-bord.png}{Tableau de bord}{fig:dashboard}

\subsubsection{R\'esultat d'une analyse}
\textit{[\`A compl\'eter]}
\TODOFigure{resultat-analyse.png}{R\'esultat d'une analyse}{fig:resultat}

\subsection{Syst\`eme de gestion de base de donn\'ees}

Pour ce projet, nous avons utilis\'e \textbf{MySQL} (MariaDB) comme syst\`eme de
gestion de base de donn\'ees relationnelle, int\'egr\'e via XAMPP. La base de donn\'ees
compte quatre tables principales~: \texttt{utilisateurs}, \texttt{analyses},
\texttt{controles} et \texttt{parametres\_score}, refl\'etant directement le
diagramme de classes pr\'esent\'e au chapitre pr\'ec\'edent. L'utilisation de XAMPP
facilite le d\'eveloppement en local en fournissant un environnement complet
incluant Apache, PHP et MySQL. Laravel interagit avec MySQL gr\^ace \`a
\textbf{Eloquent}, son ORM int\'egr\'e, qui permet de manipuler les donn\'ees de
mani\`ere simple, s\'ecuris\'ee et orient\'ee objet.

\subsection{Protocole de s\'ecurit\'e}

La s\'ecurit\'e a \'et\'e au c\oe ur de notre d\'emarche de d\'eveloppement, ce qui est
particuli\`erement coh\'erent avec la nature m\^eme du projet.

\paragraph{Authentification} L'authentification repose sur \textbf{Laravel
Sanctum} en mode SPA (authentification par cookie de session), adapt\'ee \`a la
s\'eparation frontend React / backend Laravel. Les mots de passe sont stock\'es
hach\'es (\texttt{bcrypt} via \texttt{Hash::make}), jamais en clair.

\paragraph{Autorisation} Chaque route sensible de l'API v\'erifie que
l'utilisateur authentifi\'e est bien propri\'etaire de la ressource demand\'ee (par
exemple, un utilisateur ne peut consulter le d\'etail d'une analyse que si elle
lui appartient, sauf s'il est administrateur), avec un rejet explicite (code
HTTP 403) sinon.

\paragraph{Validation des entr\'ees} Toutes les donn\'ees entrantes (URL \`a
analyser, formulaires d'inscription/connexion) sont valid\'ees c\^ot\'e serveur via
le validateur int\'egr\'e de Laravel avant tout traitement.

\subsection{Protection des donn\'ees et communication s\'ecuris\'ee}

\begin{itemize}
  \item Les injections SQL sont \'evit\'ees gr\^ace \`a l'ORM Eloquent, qui prot\`ege
    automatiquement les requ\^etes contre ces attaques.
  \item Le moteur d'analyse lui-m\^eme reste strictement \textbf{non intrusif}~: il
    se limite \`a observer les r\'eponses HTTP publiques d'un site (en-t\^etes,
    certificat) sans jamais tenter d'exploiter une faille.
  \item % TODO : une fois d\'eploy\'ee, pr\'eciser ici que les \'echanges entre le
        % frontend et l'API se font via HTTPS en production.
        \textit{[\`A compl\'eter apr\`es d\'eploiement~: confirmation du HTTPS en production]}
\end{itemize}

\section{Discussion}

\subsection{Par rapport aux objectifs fix\'es}

Les objectifs sp\'ecifiques fix\'es en introduction sont, \`a ce stade,
globalement atteints~: la plateforme authentifie ses utilisateurs (par mot de
passe ou via Google), analyse un site sur les sept crit\`eres retenus, calcule
un score et un niveau de risque, g\'en\`ere des recommandations, et conserve un
historique consultable. La principale valeur ajout\'ee par rapport aux outils
existants cit\'es en introduction (Mozilla Observatory, SSL Labs) tient moins
aux contr\^oles eux-m\^emes -- volontairement plus restreints -- qu'\`a leur
restitution~: un score unique plut\^ot que plusieurs rapports \'epars, et des
descriptions r\'edig\'ees en fran\c{c}ais courant plut\^ot qu'en jargon technique
anglophone. Ce dernier point n'\'etait pas acquis d\`es le d\'epart~: les
premi\`eres versions du moteur produisaient des constats bruts (noms
d'en-t\^etes HTTP, versions de protocole) qui se sont r\'ev\'el\'es tout aussi
incompr\'ehensibles pour un non-technicien que ceux des outils existants. La
reformulation compl\`ete des descriptions et recommandations a \'et\'e n\'ecessaire
pour r\'epondre effectivement \`a l'objectif de p\'edagogie.

\subsection{Limites}

Plusieurs limites doivent \^etre soulign\'ees. D'abord, le score refl\`ete un
\'etat \`a un instant T~: une analyse est une photographie, pas une surveillance
continue -- un site conforme au moment de l'analyse peut se d\'egrader ensuite
sans que l'utilisateur en soit inform\'e, hors nouvelle analyse manuelle.
Ensuite, les sept contr\^oles retenus, bien que repr\'esentatifs, ne constituent
pas un audit de s\'ecurit\'e exhaustif~: ils ne remplacent pas un test
d'intrusion men\'e par un professionnel, et se limitent volontairement \`a des
v\'erifications passives. La pond\'eration des crit\`eres dans le calcul du score,
bien que configurable par un administrateur, reste \'egalement une d\'ecision en
partie subjective~: deux baremes diff\'erents peuvent produire deux niveaux de
risque diff\'erents pour un m\^eme site.

Le traitement des analyses d\'ependant d'appels r\'eseau vers des serveurs
externes r\'eels, le r\'esultat peut en outre varier d'une ex\'ecution \`a
l'autre~: au cours des tests, une m\^eme URL a pu \^etre marqu\'ee \og~injoignable~\fg{}
lors d'un essai puis analys\'ee avec succ\`es quelques minutes plus tard,
simplement parce que le site cibl\'e a mis plus de huit secondes \`a r\'epondre la
premi\`ere fois. Ce comportement n'est pas un dysfonctionnement de la
plateforme, mais une contrainte inh\'erente \`a toute analyse d\'ependante du
r\'eseau, qu'il convient de garder \`a l'esprit dans l'interpr\'etation d'un
\'echec ponctuel.

Enfin, le traitement asynchrone des analyses (section~\ref{chap:programmation})
repose sur un worker de file d'attente qui doit tourner en continu~: en
environnement de d\'eveloppement local, l'oubli de d\'emarrer ce processus a \`a
plusieurs reprises laiss\'e des analyses bloqu\'ees ind\'efiniment au statut
\og~en cours~\fg, sans message d'erreur explicite pour l'utilisateur. Ce point
constitue une vigilance op\'erationnelle \`a documenter pour un futur
d\'eploiement en production, au-del\`a du seul code applicatif.

Ce chapitre a permis de pr\'esenter les principaux r\'esultats techniques du
projet ainsi qu'une lecture critique de leur port\'ee et de leurs limites.


% ===================== CONCLUSION =====================
\chapter*{Conclusion et perspectives}
\addcontentsline{toc}{chapter}{Conclusion et perspectives}
\markboth{Conclusion et perspectives}{}

\section{Conclusion}

Le d\'eveloppement de la plateforme CyberGuard B\'enin a \'et\'e l'occasion de mettre
en pratique l'ensemble du cycle de vie d'un projet informatique~: clarification
du besoin, mod\'elisation UML, conception de la base de donn\'ees, puis
impl\'ementation technique. Bas\'ee sur une architecture combinant une API Laravel
et une interface React, la plateforme automatise un diagnostic de s\'ecurit\'e
jusque-l\`a r\'eserv\'e \`a des outils techniques et anglophones, en le rendant
accessible \`a un public non-technicien gr\^ace \`a un score, un niveau de risque et
des recommandations en fran\c{c}ais.

Au terme du stage, la plateforme livr\'ee couvre l'ensemble des cas d'utilisation
identifi\'es en phase de conception~: authentification (par mot de passe et via
Google), lancement et traitement asynchrone des analyses (sept contr\^oles non
intrusifs), consultation des rapports avec t\'el\'echargement PDF, tableau de
bord, notifications in-app de fin d'analyse, ainsi qu'un back-office
d'administration complet (gestion des comptes, mod\'eration des analyses,
param\'etrage du bar\`eme de scoring). Le principal point rest\'e en dehors du
p\'erim\`etre du stage est le d\'eploiement en production, la plateforme n'ayant
\'et\'e ex\'ecut\'ee qu'en environnement de d\'eveloppement local.

\section{Perspectives}

Parmi les perspectives d'\'evolution envisag\'ees~:
\begin{itemize}
  \item le d\'eploiement en production et l'ouverture de la plateforme au
    public, avec confirmation du HTTPS de bout en bout~;
  \item la v\'erification de l'adresse email \`a l'inscription (compte cr\'e\'e par
    formulaire classique)~;
  \item le suivi automatique p\'eriodique~: relancer r\'eguli\`erement l'analyse
    d'un site d\'ej\`a enregistr\'e et alerter par email en cas de r\'egression du
    score, plut\^ot que de d\'ependre uniquement d'une analyse manuelle~;
  \item un rapport public partageable (lien en lecture seule) et un badge de
    score int\'egrable sur le site audit\'e~;
  \item l'ajout de v\'erifications suppl\'ementaires (CSP plus fine, en-t\^etes
    additionnels, historique des CVE connues li\'ees aux technologies
    d\'etect\'ees)~;
  \item une couverture de tests automatis\'es \'etendue au frontend, en
    compl\'ement de la suite PHPUnit existante c\^ot\'e backend.
\end{itemize}

En d\'efinitive, CyberGuard B\'enin se veut une contribution concr\`ete \`a la
d\'emocratisation de la cybers\'ecurit\'e au B\'enin, en r\'eduisant la surface
d'attaque des sites web mal s\'ecuris\'es et en rendant la premi\`ere \'etape d'un
audit de s\'ecurit\'e accessible \`a tous.


% ===================== REFERENCES =====================
\chapter*{R\'ef\'erences bibliographiques}
\addcontentsline{toc}{chapter}{R\'ef\'erences bibliographiques}
\markboth{R\'ef\'erences bibliographiques}{}

\begin{enumerate}
  \item Laravel -- The PHP Framework For Web Artisans. Consult\'e en 2026 sur
    \url{https://laravel.com/docs}
  \item Laravel Sanctum -- SPA Authentication. Consult\'e en 2026 sur
    \url{https://laravel.com/docs/sanctum}
  \item React Documentation. Consult\'e en 2026 sur \url{https://react.dev}
  \item Mozilla Observatory. Consult\'e en 2026 sur \url{https://developer.mozilla.org/en-US/observatory}
  \item Code de statut HTTP -- MDN Web Docs. Consult\'e en 2026 sur
    \url{https://developer.mozilla.org/fr/docs/Web/HTTP/Reference/Status}
  % TODO : ajouter toute autre source r\'eellement consult\'ee pendant le stage.
\end{enumerate}


% ===================== ANNEXES =====================
\chapter*{Annexes}
\addcontentsline{toc}{chapter}{Annexes}
\markboth{Annexes}{}

\section*{1. Extrait du code pour le lancement d'une analyse}
\addcontentsline{toc}{section}{1. Extrait du code pour le lancement d'une analyse}

\begin{lstlisting}[language=PHP]
/**
 * Cas d'utilisation : Lancer une analyse -> declenche "Traiter l'analyse".
 */
public function store(Request $request)
{
    $data = $request->validate([
        'url' => ['required', 'string', 'max:255', 'url'],
    ]);

    $analyse = $request->user()->analyses()->create([
        'url' => $data['url'],
        'date_analyse' => now(),
        'statut' => 'en_cours',
    ]);

    try {
        $this->service->executer($analyse);
    } catch (\Throwable $e) {
        Log::error('Echec du traitement de l\'analyse #'.$analyse->id, ['exception' => $e]);
        $analyse->update(['statut' => 'echec']);
    }

    return response()->json($analyse->fresh('resultatsVerification'), 201);
}
\end{lstlisting}

\section*{2. Extrait du moteur de v\'erification (calcul du score)}
\addcontentsline{toc}{section}{2. Extrait du moteur de v\'erification (calcul du score)}

\begin{lstlisting}[language=PHP]
private function calculerScore(array $resultats, ParametreScore $bareme): array
{
    $poidsParType = [
        'https' => $bareme->poids_https / 2,
        'certificat_ssl' => $bareme->poids_https / 2,
        'en_tetes_securite' => $bareme->poids_headers,
        'cookies' => $bareme->poids_cookies,
        'vulnerabilites' => $bareme->poids_vulnerabilites,
    ];

    $score = 0;
    foreach ($resultats as $resultat) {
        if ($resultat['statut'] === 'conforme') {
            $score += $poidsParType[$resultat['type']] ?? 0;
        }
    }

    $score = (int) round($score);

    $niveauRisque = match (true) {
        $score >= $bareme->seuil_bon => 'bon',
        $score >= $bareme->seuil_moyen => 'moyen',
        default => 'faible',
    };

    return [$score, $niveauRisque];
}
\end{lstlisting}

% TODO : ajouter d'autres extraits pertinents au fil du d\'eveloppement
% (migrations, composants React, etc.).


% ===================== TABLE DES MATIERES DETAILLEE =====================
\tableofcontents

\end{document}